\chapter{Conclusões e Trabalhos Futuros} \label{cap:conclusoes-e-trabalhos-futuros}

    A competição de robôs solucionadores de labirintos, os \textit{micromouses}, já é, há algum tempo, uma realidade consolidade dentro do cenário do estudo de robótica e controle em especial no cenário nacional haja vista seu caráter multidisciplinar no âmbito da Engenharia. E, dessa forma, o desenvolvimento de um robô para participação de competições do tipo é um trabalho relativamente complexo e muitas vezes realizado em equipe. Nesse contexto, o presente trabalho buscou abordar uma das etapas iniciais nesse desenvolviment, isto é, a etapa de controle das velocidades do robô.

    Tendo em vista a metodologia desenvolvida no que se refere à comunicação sem fio, pode-se concluir que esta demonstrou-se adequada dentro do cenário do Laboratório de Controle e Automação do grupo \gls{GRASI} do curso de Engenharia Elétrica da \gls{UFPI}, uma vez que as coletas de dados e a identificação evidenciaram um correto funcionamento da telemetria. Além disso, a comunicação sem fio mostrou-se compatível com a natureza móvel do robô, permitindo a execução de testes sem a utilização de conexões físicas que pudessem restringir seu deslocamento. Dessa forma, foi possível realizar experimentos em superfícies e condições de operação mais próximas daquelas encontradas em ambientes de competição.

    Já no referente à metodologia adotada para o projeto de controladores é possível observar, para os controladores dos motores, que, apesar de ruídos de medição que possam aparecer no processo ou de não-linearidades não modeladas, os motores mantiveram suas velocidades estáveis e controaldas, seguindo as diversas referências testadas. Quanto aos controladores supervisórios os resultados apresentados levam a crer que, com o desacoplamento inserido devido à utilização de malhas internas, o comportamento do robô também se manterá estável se realizados testes em condições reais.

    % Assim, conclui-se que o robô micromouse, apesar da sua complexidade, pode ser, de maneira relativamente fácil, adequadamente controlado quando em condições adequadas de testagem e dispondo de meios confiáveis para a transmissão das informações obtidas.

    Diante dos resultados apresentados, conclui-se que os objetivos relacionados à aquisição de dados, à identificação dos sistemas e ao projeto dos controladores de velocidade foram alcançados nas condições experimentais consideradas. A estrutura desenvolvida possibilitou o controle individual dos motores e estabeleceu uma base para o controle das variáveis globais de movimento do \textit{micromouse}. 
    
    Assim, o trabalho apresenta-se como um passo importante para as etapas posteriores de desenvolvimento da plataforma, especialmente aquelas relacionadas à navegação autônoma e à resolução de labirintos.

    Para trabalhos futuros, e de posse da estrutura já desenvolvida, sugere-se abordar o desenvolvimento de algorítmos de rastreio de trajetórias com e sem o uso de sistemas externos de visão, tais quais câmeras. Tais algorítmos serviriam de base ainda, para o entendimento e elaboração de algoritmos de resolução de labirintos como o algoritmo \textit{Flood Fill}.

    







